\section{Research Objectives} \label{section:research_objectives}

\label{para:mlp_interface_aims} % ---------- %

This project aims to evaluate and enhance the use of machine-learned interatomic potentials (MLPs) for predicting the
relative stability of material interfaces.

By replacing expensive density functional theory (DFT) calculations in the early stages of interface screening, MLPs
offer a route to scalable evaluation across diverse chemical systems and stacking configurations.

\label{para:interface_config_space} % ---------- %

The configurational landscape of interface structures is combinatorially large.

Even for lattice-matched systems, variations in Miller planes, interfacial shifts, and terminations can lead to
significant changes in formation energy and electronic response.

Traditional DFT methods, while accurate, are computationally prohibitive when applied to hundreds or thousands of such
configurations.

\label{para:mlp_pipeline_integration} % ---------- %

This work builds on the integration of MLPs, particularly the MACE model, into high-throughput interface generation and
evaluation pipelines.

It focuses not only on energy prediction, but also on favourability ranking, robustness to system-specific errors, and
hybrid workflows where MLPs are used as a pre-screening stage ahead of final DFT validation.

\label{para:core_objectives} % ---------- %

\paragraph{The research addresses the following core objectives:}

\begin{enumerate}

    \item \textbf{Benchmark MLP-predicted rankings against DFT across diverse interfaces.}
    Using MACE-evaluated relaxations and energies, this study quantifies how well MLPs preserve the DFT ordering of
    interface favourability across multiple group-IV homostructures and heterostructures.

    \item \textbf{Diagnose system-specific discrepancies and explore delta-learning corrections.}
    Particular attention is given to cases, such as pure-Si systems, where MLP rank fidelity degrades.

    Chapter~\ref{chapter:next_steps} proposes training delta models to correct these systematic energy errors.

    \item \textbf{Develop structure-based predictors of interface energy.}
    Surface and slab-level features (e.g. orientation, element type, symmetry, strain) will be used to train models
    capable of predicting favourability prior to full interface assembly, reducing reliance on supercell enumeration.

    \item \textbf{Integrate MLPs and predictors into automated workflows.}
    Tools such as ARTEMIS and RAFFLE will be extended to incorporate MLP-based scoring or filtering, enabling efficient
    identification of favourable configurations before expensive DFT refinement.

\item \end{enumerate}

\label{para:scalable_screening_implications} % ---------- %

In the context of this report, the MLP‐driven workflow enables rapid, scalable screening of interface configurations and
informs the design of functional heterostructures by ranking candidates according to atomic‐level energetics.

By focusing on regimes, such as mismatched or metastable reconstructions, where traditional coherence rules or
bulk‐derived heuristics break down, this approach hopes to uncover non‐intuitive structures that would otherwise be overlooked.